\pdfoutput=1
\documentclass[11pt]{amsart}
\usepackage{amsmath,amssymb,amsthm}
\usepackage[margin=1.2in]{geometry}
\usepackage{booktabs,adjustbox,xcolor}
\usepackage[colorlinks=true,linkcolor=blue,citecolor=blue,urlcolor=blue]{hyperref}
\hypersetup{pdfauthor={Bernd Johannes Wuebben},
  pdftitle={Small eigenvalues on degenerating curves and the intersection form}}
\setlength{\emergencystretch}{2em}
\newtheorem{theorem}{Theorem}[section]
\newtheorem{proposition}[theorem]{Proposition}
\newtheorem{lemma}[theorem]{Lemma}
\newtheorem{corollary}[theorem]{Corollary}
\theoremstyle{definition}
\newtheorem{definition}[theorem]{Definition}
\newtheorem{example}[theorem]{Example}
\theoremstyle{remark}
\newtheorem{remark}[theorem]{Remark}
\newcommand{\ZZ}{\mathbb Z}
\newcommand{\RR}{\mathbb R}
\newcommand{\CC}{\mathbb C}
\newcommand{\PP}{\mathbb P}
\newcommand{\Sing}{\operatorname{Sing}}
\newcommand{\area}{\operatorname{Area}}
\newcommand{\dc}{d^c}
\title[Small eigenvalues on degenerating curves]{Small eigenvalues on degenerating curves\\and the intersection form}
\author{Bernd Johannes Wuebben}
\subjclass[2020]{58J50 (primary); 14D06, 32G20, 05C50 (secondary)}
\keywords{Laplacian, small eigenvalue, degeneration of Riemann surfaces,
intersection form, dual graph, plane curve singularity}
\date{November 10, 2026\\This version: November 10, 2026}
\begin{document}
\begin{abstract}
Let a holomorphic family of compact Riemann surfaces over a disc, with
smooth total space, degenerate to a reduced curve with $N$ irreducible
components and arbitrary singularities, and let the fibres carry the
metrics induced by a Riemannian metric on the total space that is
Hermitian near the singular points of the central fibre. We prove that
the $N-1$ eigenvalues of the Laplace--Beltrami operator on functions that
tend to zero, multiplied by $\log(1/|s|)$, where $s$ is the parameter,
converge to $2\pi$ times the eigenvalues
of the negative intersection matrix of the components relative to the
diagonal matrix of their areas, and that the remaining positive
eigenvalues stay bounded away from zero. Equivalently, the small
eigenvalues are asymptotic to those of a finite weighted graph on the
components whose edge weights are intersection numbers, in analogy with
Burger's theorem for hyperbolic surfaces. This answers questions of Dai
and Yoshikawa on the existence and geometric meaning of these limits, and
evaluates the constant in their asymptotic formula for the product of the
small eigenvalues as a weighted count of spanning trees. For K\"ahler
metrics the limits can also be derived from recent work of Cao on the
Archimedean height pairing combined with the eigenvalue bounds of Dai and
Yoshikawa; the proof given here is direct. It uses the logarithms of local equations of the components
as explicit harmonic functions on the smooth fibres, whose Dirichlet
energies are computed by winding numbers.
\end{abstract}
\maketitle

\section{Introduction}\label{sec:intro}

\subsection{The result}
Let $X$ be a connected complex surface, $\mathbb D\subset\CC$ a disc centred
at the origin, and
\[
 \pi:X\longrightarrow\mathbb D
\]
a proper surjective holomorphic map with connected fibres $X_s=\pi^{-1}(s)$,
which is a submersion over $\mathbb D\setminus\{0\}$. We assume that the
central fibre $X_0$ is a reduced divisor; equivalently, $\pi$ has finitely
many critical points, and these are the singular points of $X_0$: a
reduced local equation of a curve has vanishing differential exactly at
its singular points, and a multiple component would be critical along a
curve. The fibre $X_0$ is connected, as a limit of connected fibres of a
proper map. Let
$C_1,\ldots,C_N$ be the irreducible components of $X_0$. No further
condition is imposed on the singularities of $X_0$.

Let $g$ be a smooth Riemannian metric on $X$ that is Hermitian on a
neighbourhood of $\Sing X_0$, and let $g_s$ be its restriction to $X_s$ for $s\ne0$, and $g_0$ its
restriction to $X_0\setminus\Sing X_0$.
For $s\ne0$ let
\[
 0=\lambda_0(s)<\lambda_1(s)\le\lambda_2(s)\le\cdots
\]
be the eigenvalues of the Laplacian $\Delta_s=d^*d$ of $(X_s,g_s)$ acting
on functions. Let $A_i$ be the area of $C_i\setminus\Sing X_0$ with respect
to $g$.

On $\RR^N$ consider the two quadratic forms
\begin{equation}\label{eq:forms}
 q(F)=\sum_{i<j}(C_i\cdot C_j)\,(F_i-F_j)^2,\qquad
 m(F)=\sum_{i}A_iF_i^2,
\end{equation}
where $C_i\cdot C_j$ is the intersection number in $X$. Since
$\sum_iC_i=X_0$ is a fibre, $C_i\cdot X_0=0$, and therefore
\begin{equation}\label{eq:q-intersection}
 q(F)=-\Bigl(\sum_iF_iC_i\Bigr)^2 .
\end{equation}
Thus $q$ is the negative of the intersection form on the span of the
components. It is positive semidefinite, and its kernel consists of the
constant vectors because $X_0$ is connected; this is Zariski's lemma
\cite[Chapter~III, Lemma~8.2]{bhpv}, and in the present form it is also
evident from \eqref{eq:forms}. Let
\[
 0=\mu_0<\mu_1\le\cdots\le\mu_{N-1}
\]
be the eigenvalues of $q$ relative to $m$, that is, the roots of
$\det(Q-\mu\operatorname{diag}(A_i))=0$, where $Q$ is the matrix of $q$.

\begin{theorem}\label{thm:main}
With the preceding notation,
\begin{equation}\label{eq:main}
 \lim_{s\to0}\,\log(1/|s|)\,\lambda_k(s)=2\pi\mu_k
 \qquad(1\le k\le N-1),
\end{equation}
and $\liminf_{s\to0}\lambda_N(s)>0$.
\end{theorem}

The limit depends on the metric only through the areas of the components,
and on the singularities of $X_0$ only through the local intersection
multiplicities of distinct components. For a finite graph with
weights (conductances) $c_e>0$ on the edges and masses $m_v>0$ at the
vertices, we call the eigenvalues of the form
$\sum_ec_e(F_{e^+}-F_{e^-})^2$ relative to $\sum_vm_vF_v^2$ the eigenvalues
of the weighted graph. If $X_0$ has only ordinary nodes,
the matrix of $q$ is the Laplacian of the dual graph of $X_0$, and
\eqref{eq:main} says that the small eigenvalues are asymptotic to those of
the dual graph with conductance $2\pi/\log(1/|s|)$ on each edge and mass
$A_i$ at each vertex. A point where two components are simply tangent
counts as two edges, and an ordinary $d$-fold point through which $d$
components pass contributes a complete graph on those components
(Section~\ref{sec:examples}).

\begin{corollary}\label{cor:two}
If $N=2$, then
\[
 \lambda_1(s)=\frac{2\pi\,(C_1\cdot C_2)\,(A_1+A_2)}{A_1A_2}\,
 \frac{1+o(1)}{\log(1/|s|)} .
\]
\end{corollary}

\begin{corollary}\label{cor:product}
Let $\kappa$ be the common value of the principal cofactors of the matrix
of $q$; by the matrix-tree theorem \cite{chaiken} it is the sum over the spanning trees
$\mathcal T$ of the complete graph on $\{1,\ldots,N\}$ of
$\prod_{\{i,j\}\in\mathcal T}(C_i\cdot C_j)$. Then
\[
 \prod_{k=1}^{N-1}\lambda_k(s)
 =\Bigl(\frac{2\pi}{\log(1/|s|)}\Bigr)^{N-1}
 \Bigl(\frac{\kappa\,\sum_iA_i}{\prod_iA_i}+o(1)\Bigr).
\]
\end{corollary}

\subsection{Relation to earlier work}
For compact hyperbolic surfaces, Schoen, Wolpert and Yau \cite{swy}
bounded the small eigenvalues above and below by sums of lengths of short
geodesics, and Burger \cite[Theorem~1.1]{burger} proved that they are
asymptotic to the eigenvalues of a weighted graph: the vertices are the
components of the complement of the short geodesics, with masses
$2\gamma_v-2+n_v$, where $\gamma_v$ is the genus and $n_v$ the number of
boundary geodesics and cusps of the component, and the edges are the
geodesics, with weights their lengths; the ratio of each small eigenvalue
of the surface to the corresponding eigenvalue of the graph tends to
$1/(2\pi^2)$.

For metrics induced from the total space, Dai and Yoshikawa
\cite{dai-yoshikawa} work in the setting above, with $X$ a complex
surface carrying a positive line bundle, $g$ K\"ahler and $X_0$ reduced
and reducible. They prove that
$\lambda_1,\ldots,\lambda_{N-1}$ are bounded above and below by constant
multiples of $1/\log(1/|s|)$ \cite[Theorem~0.2]{dai-yoshikawa}, and that
$(\log(1/|s|))^{N-1}\prod_{k=1}^{N-1}\lambda_k(s)$ has a positive limit
\cite[Theorem~0.3]{dai-yoshikawa}; in particular
$\lambda_1(s)\log(1/|s|)$ converges when $N=2$. The proof compares ratios of analytic torsions with determinants of period
integrals. The constant is not evaluated. They ask whether
$\lambda_k(s)\log(1/|s|)$ converges when $N>2$, for a geometric expression
of the limit, and whether there is a finite weighted graph whose
eigenvalues are asymptotic to the small eigenvalues
\cite[Problems~9.2 and~9.4]{dai-yoshikawa}. Theorem~\ref{thm:main}
answers these questions in the setting of that paper: the graph has
vertices the components, masses $A_i$ and conductances
$2\pi(C_i\cdot C_j)/\log(1/|s|)$, and Corollary~\ref{cor:product}
evaluates the constant in their product formula. Their Laplacian is
$\bar\partial^*\bar\partial=\tfrac12d^*d$, so in their normalization
$2\pi$ is replaced by~$\pi$.

Cao \cite{cao-kahler} extends the two-sided bounds to the scalar Laplacian
on degenerations of K\"ahler manifolds of any dimension, and conjectures
that the rescaled small eigenvalues converge to those of a weighted dual
graph, without specifying the edge weights
\cite[Theorem~0.4 and Section~3.C]{cao-kahler}. Theorem~\ref{thm:main}
confirms this for degenerations of curves with reduced central fibre and
identifies the edge weights as the intersection numbers $C_i\cdot C_j$. In \cite{cao-pairing} he
studies, for degenerating curves, the Archimedean height pairing of two-forms, defined by
solving $dd^c\varphi=-\alpha$ on the fibres (with $d^c$ normalized as in
that paper; see Remark~\ref{rem:cao}), and proves that its
logarithmic coefficient is given by the Moore--Penrose inverse of the
intersection matrix of the components
\cite[Theorem~0.2]{cao-pairing}. That result concerns the Green operator on the small eigenspaces, and
Theorem~\ref{thm:main} the eigenvalues themselves. For K\"ahler metrics,
which include the setting of \cite{dai-yoshikawa}, the limits
\eqref{eq:main} can also be derived from \cite{cao-pairing} and the bounds
of \cite{dai-yoshikawa} by a finite-dimensional argument that is carried
out in neither paper; we sketch it in Remark~\ref{rem:cao}. What is added
here is a direct proof that does not use the Green operator or analytic
torsion, the extension to metrics that are only Hermitian near
$\Sing X_0$, and the explicit identification of the graph and of the
constants.

Ji and Wentworth \cite{ji-wentworth} prove spectral convergence for
plumbing families with the Bergman metric and, for a separating node,
suggest a heuristic for the rate of $\lambda_1$
\cite[Remark~5.10]{ji-wentworth}; see Remark~\ref{rem:jw}.

\subsection{Method}
Near a singular point $p$ of $X_0$ write $\pi=\prod_if_i$, where $f_i$ is
a local equation of $C_i$. On a smooth fibre the functions $\log|f_i|$ are
harmonic and sum to $\log|s|$. For $F\in\RR^N$ the function
\[
 h_F=\frac1{\log|s|}\sum_iF_i\log|f_i|
\]
is harmonic, takes values between $\min F_i$ and $\max F_i$, and is close
to $F_i$ on the part of the fibre near $C_i$ and away from $p$. By
Stokes' theorem its Dirichlet energy is a sum of boundary terms, each the
product of a boundary value $F_i$ and a winding number of some $f_j$ along
a boundary circle. These winding numbers are local intersection
multiplicities, and the energy is $2\pi/\log(1/|s|)$ times the local
contribution to $q(F)$, up to an error of order $|F|^2/\log(1/|s|)^2$
(Lemma~\ref{lem:energy}). Used as test functions the $h_F$ give the upper
bound; they play the role of the capacitary potentials of the
neighbourhoods of the singular points in the fibre. For the lower bound we pair the differential of an arbitrary
function of small energy with $\dc h_F$, where $\dc=i(\bar\partial-\partial)$; the Cauchy--Schwarz inequality
then bounds its energy below by the same quantity
(Lemma~\ref{lem:capacity}). No resolution of the singularities of $X_0$
and no semistable reduction is used.

\subsection{Conventions}
We write $\dc=i(\bar\partial-\partial)$. On a Riemann surface with a
conformal metric, $\dc u=*du$ for real functions $u$, so that
$du\wedge\dc u=|du|^2\,dA$, and a function $u$ is harmonic if and only if
$d\dc u=0$. If $f$ is holomorphic and nowhere zero, $\dc\log|f|=d\arg f$.
The Dirichlet energy of $u$ over a domain $M$ is
$E_M(u)=\int_M|du|^2\,dA$, with associated bilinear form $E_M(u,v)$;
we write $E=E_{X_s}$. Norms without subscript are $L^2$ norms on
$(X_s,g_s)$. The letter $C$ denotes constants that depend only on
$(X,\pi,g)$, and $c(r)$ denotes constants that depend in addition on a
radius $r$; neither depends on $s$, and their values may change from line
to line. For $F\in\RR^N$ we put $|F|=\max_i|F_i|$.

\section{Local structure near the singular points}\label{sec:local}

\subsection{Coordinates and local equations}
Let $\Sigma=\Sing X_0$. For each $p\in\Sigma$ choose holomorphic
coordinates $z=(z_1,z_2)$ centred at $p$ on an open set $U_p$ on which $g$
is Hermitian, such that neither $\pi(z_1,0)$ nor $\pi(0,z_2)$ vanishes
identically. Write $B_p(r)=\{|z|<r\}$ and $S_p(r)=\partial B_p(r)$.
Since $X_0$ is reduced, the germ of $\pi$ at $p$ is a product of pairwise
non-associated irreducible germs, one for each local branch $\beta$ of
$X_0$ at $p$. Collecting the branches that lie on the same component, and
absorbing the unit into one factor, we obtain holomorphic functions
$f_{1},\ldots,f_{N}$ on $U_p$, after shrinking it, with
\begin{equation}\label{eq:factor}
 \pi=\prod_{i=1}^Nf_i\quad\text{on }U_p,\qquad
 \{f_i=0\}=C_i\cap U_p ,
\end{equation}
where $f_i$ vanishes to first order along $C_i\cap U_p$ and
$f_i\equiv1$ if $p\notin C_i$. Shrinking $U_p$ again, $|f_i|<1$ whenever
$p\in C_i$. For a branch $\beta$ let $i(\beta)$ be the index of the
component containing it, and for $j\ne i(\beta)$ let $(\beta\cdot C_j)_p$
be the intersection multiplicity of $\beta$ with $C_j$ at $p$, that is,
the order of vanishing at $p$ of $f_j$ along the normalization of $\beta$.
Then
\begin{equation}\label{eq:local-intersection}
 \sum_{\beta\subset C_i}(\beta\cdot C_j)_p=(C_i\cdot C_j)_p,\qquad
 \sum_{p\in\Sigma}(C_i\cdot C_j)_p=C_i\cdot C_j\qquad(i\ne j).
\end{equation}

\subsection{Radii}
Each branch $\beta$ of $X_0$ at $p$ has an injective holomorphic
parametrization $\tau\mapsto\phi(\tau)$ on a disc, with $\phi(0)=p$ and
$|\phi(\tau)|^2=c|\tau|^{2n}(1+O(|\tau|))$ for some $c>0$ and $n\ge1$, so
that $|\phi|$ is strictly increasing along rays for small $|\tau|$. Hence
there is $\rho_0>0$ such that for every $p\in\Sigma$ the ball
$\overline{B_p(4\rho_0)}$ lies in $U_p$, these closed balls are pairwise
disjoint, and for $0<r\le4\rho_0$ the sphere $S_p(r)$ meets $X_0$
transversally, each branch $\beta$ meeting $B_p(r)$ in the image of a
disc $\{|\phi|<r\}$ that is star-shaped about $0$, whose boundary circle
$\gamma_\beta^0(r)=\beta\cap S_p(r)$ therefore winds once around $p$ in
the normalization of $\beta$; compare the conic structure theorem
\cite[Section~2]{milnor}. For $0<r\le2\rho_0$ put
\[
 K_r=X\setminus\bigcup_{p\in\Sigma}B_p(r),\qquad
 \Omega_i^0(r)=C_i\cap K_r .
\]
Each $\Omega_i^0(r)$ is a compact connected surface with boundary: its
preimage in the normalization of $C_i$ is the complement of finitely many
discs. The sets $\Omega^0_i(r)$, $1\le i\le N$, are pairwise disjoint.

\begin{lemma}\label{lem:thick}
For every $r\in(0,2\rho_0]$ there are $\delta(r)>0$ and a smooth map
\[
 \Phi:(X_0\cap K_r)\times\{|s|<\delta(r)\}\longrightarrow K_r,
\]
a diffeomorphism onto $\pi^{-1}(\{|s|<\delta(r)\})\cap K_r$, with
$\pi\circ\Phi(x,s)=s$ and $\Phi(x,0)=x$. For $0<|s|<\delta(r)$:
\begin{enumerate}
\item $X_s\cap K_r$ is the disjoint union of the connected surfaces
$\Omega_i(r,s)=\Phi(\Omega^0_i(r),s)$, and the metrics
$\Phi(\cdot,s)^*g_s$ converge to $g_0$ in every $C^k$ norm on
$X_0\cap K_r$ as $s\to0$;
\item $X_s$ meets each $S_p(r)$ transversally, in a disjoint union of
circles $\gamma_\beta(r,s)$, one for each branch $\beta$ at $p$,
with $\gamma_\beta(r,s)\to\gamma^0_\beta(r)$ in $C^1$ as $s\to0$;
\item $M_p(r,s)=X_s\cap\overline{B_p(r)}$ is a compact surface with
boundary $\bigcup_\beta\gamma_\beta(r,s)$.
\end{enumerate}
\end{lemma}
\begin{proof}
The map $\pi$ has no critical points on $X_0\cap K_r$, and its restriction
to the boundary spheres has none on $X_0\cap\partial K_r$, by
transversality. The set of points of $K_r$ at which $\pi$ or
$\pi|_{\partial K_r}$ fails to be a submersion is closed and disjoint from
$X_0$; since $\pi$ is proper, its image misses a disc
$\{|s|<\delta(r)\}$. Over this disc the real coordinate vector fields of
$\mathbb D$ have lifts to $K_r$ that are tangent to $\partial K_r$: such
lifts exist locally, at boundary points because $\pi|_{\partial K_r}$ is a
submersion, and the two conditions are preserved by a partition of unity.
Composing the flows of the two lifts gives $\Phi$, as in Ehresmann's
theorem. Assertion (1) follows because $\Phi$ is smooth and
$\Phi(\cdot,0)$ is the identity. Since $\pi|_{S_p(r)}$ is a submersion
along $X_s\cap S_p(r)$, the fibre $X_s$ is transverse to $S_p(r)$, and
$X_s\cap S_p(r)=\Phi(X_0\cap S_p(r),s)$; this gives (2) and (3).
\end{proof}

We always take $|s|$ smaller than $\delta$ of every radius involved, without further mention. The maps $\Phi$ for different radii
need not agree, but the sets they define are compatible for small $s$:
$\Omega_i(r,s)$ is the unique connected component of $X_s\cap K_r$
contained in a prescribed small neighbourhood of $\Omega^0_i(r)$ in $X$,
because $\Phi^{-1}$ is continuous and the $\Omega^0_j(r)$ are compact and
pairwise disjoint. Consequently, for $r'<r$ and small $s$,
$\Omega_i(r,s)\subset\Omega_i(r',s)$, the circle $\gamma_\beta(r,s)$
lies in the interior of $\Omega_{i(\beta)}(r',s)$ and is a boundary
component of $\Omega_{i(\beta)}(r,s)$, and both lie in the set $W_\beta$
of Lemma~\ref{lem:collar} below.

\begin{lemma}\label{lem:area}
After decreasing $\rho_0$ there is $C$ such that
$\area(X_s\cap B_p(r))\le Cr^2$ for all $p\in\Sigma$, all
$r\in(0,4\rho_0]$ and all sufficiently small $s$, including $s=0$.
Consequently
\begin{equation}\label{eq:area-limit}
 A_i-Cr^2\le\lim_{s\to0}\area\Omega_i(r,s)=\area\Omega^0_i(r)\le A_i .
\end{equation}
\end{lemma}
\begin{proof}
Let $n_2$ be the order of vanishing of $\pi(0,z_2)$ at $z_2=0$. Choose
$\rho''>0$ with $\pi(0,z_2)\ne0$ for $0<|z_2|\le\rho''$. By Rouch\'e's
theorem there are $\rho'>0$ and $\delta'>0$ such that for $|z_1|\le\rho'$
and $|s|\le\delta'$ the function $z_2\mapsto\pi(z_1,z_2)-s$ has exactly
$n_2$ zeros in $\{|z_2|<\rho''\}$, counted with multiplicity. Thus, if
$4\rho_0\le\min(\rho',\rho'')$, the projection of $X_s\cap B_p(r)$ to the
$z_1$-axis has image in the disc of radius $r$ and at most $n_2$ points in
each fibre, so
$\int_{X_s\cap B_p(r)}\tfrac i2dz_1\wedge d\bar z_1\le n_2\pi r^2$.
The same holds with the two coordinates exchanged. The sum of the two
integrals is the area of the complex curve $X_s\cap B_p(r)$ for the flat
metric of the chart, and $g$ is bounded by a constant multiple of that
metric on $B_p(4\rho_0)$. The convergence in \eqref{eq:area-limit} is
Lemma~\ref{lem:thick}(1), and the inequalities follow from the first
assertion at $s=0$.
\end{proof}

\subsection{The harmonic functions}
Fix $p\in\Sigma$. For $s\ne0$ the functions $f_i$ of \eqref{eq:factor} are
holomorphic and nowhere zero on $X_s\cap U_p$, and
$\sum_i\log|f_i|=\log|s|$ there. For $F\in\RR^N$ define
\begin{equation}\label{eq:hF}
 h_F=\frac1{\log|s|}\sum_{i=1}^NF_i\log|f_i|
 \qquad\text{on }X_s\cap U_p .
\end{equation}
It is harmonic and linear in $F$. The numbers
$t_i=\log|f_i|/\log|s|$ are nonnegative for $|s|<1$ and sum to one, so
\begin{equation}\label{eq:hF-bounds}
 \min_iF_i\le h_F\le\max_iF_i .
\end{equation}

\begin{lemma}\label{lem:collar}
Let $0<r\le\rho_0$. For every branch $\beta$ at $p$, with $i=i(\beta)$,
there is a neighbourhood $W_\beta$ in $X$ of
$\beta\cap(\overline{B_p(4\rho_0)}\setminus B_p(r/2))$ such that, for all
sufficiently small $s\ne0$, the set
$X_s\cap(\overline{B_p(4\rho_0)}\setminus B_p(r/2))$ is contained in the
disjoint union of the $W_\beta$, and on $X_s\cap W_\beta$
\begin{equation}\label{eq:collar}
 h_F-F_i=\frac1{\log|s|}\sum_{j\ne i}(F_j-F_i)\log|f_j|,\qquad
 |h_F-F_i|+|dh_F|_{g_s}\le\frac{c(r)\,|F|}{\log(1/|s|)} .
\end{equation}
\end{lemma}
\begin{proof}
For $j\ne i$ the function $f_j$ vanishes only on $C_j$, which meets
$\beta$ in $B_p(4\rho_0)$ only at $p$. Hence $\log|f_j|$ is smooth and
bounded with bounded differential on a neighbourhood $W_\beta$ of the
compact set $\beta\cap(\overline{B_p(4\rho_0)}\setminus B_p(r/2))$, and
these neighbourhoods may be taken pairwise disjoint. By
Lemma~\ref{lem:thick} the fibre $X_s$ lies in their union over the given
region for small $s$. The identity in \eqref{eq:collar} follows from
$\log|f_i|=\log|s|-\sum_{j\ne i}\log|f_j|$, and the estimate from the
identity.
\end{proof}

Define the local forms
\begin{equation}\label{eq:qp}
 q_p(F,G)=\sum_{i<j}(C_i\cdot C_j)_p\,(F_i-F_j)(G_i-G_j),
 \qquad q=\sum_{p\in\Sigma}q_p ,
\end{equation}
the second identity being \eqref{eq:local-intersection}; here $q(F,G)$ is
the bilinear form of \eqref{eq:forms}.

\begin{lemma}[Energy identity]\label{lem:energy}
Let $0<r\le\rho_0$ and $M=M_p(r,s)$. For $F,G\in\RR^N$,
\begin{equation}\label{eq:flux}
 \sum_\beta F_{i(\beta)}\int_{\gamma_\beta(r,s)}\dc h_G
 =\frac{2\pi}{\log(1/|s|)}\,q_p(F,G),
\end{equation}
where the circles carry the orientation induced from $M$, and
\begin{equation}\label{eq:energy}
 \Bigl|E_M(h_F,h_G)-\frac{2\pi}{\log(1/|s|)}\,q_p(F,G)\Bigr|
 \le\frac{c(r)\,|F|\,|G|}{\log(1/|s|)^2} .
\end{equation}
\end{lemma}
\begin{proof}
Let $\beta$ be a branch with $i=i(\beta)$, and let $j\ne i$. On
$W_\beta$ the form $\dc\log|f_j|=d\arg f_j$ is closed, and by
Lemma~\ref{lem:thick}(2) the circle $\gamma_\beta(r,s)$ is homotopic in
$W_\beta$ to $\gamma^0_\beta(r)$, with matching orientations. The latter
is the boundary of the disc $\beta\cap B_p(r)$ and winds once around $p$
in the normalization of $\beta$, so the argument principle gives
\[
 \int_{\gamma_\beta(r,s)}\dc\log|f_j|=2\pi\,(\beta\cdot C_j)_p
 \qquad(j\ne i).
\]
Since $\log|f_i|=\log|s|-\sum_{j\ne i}\log|f_j|$ on $X_s$,
\[
 \int_{\gamma_\beta(r,s)}\dc h_G
 =\frac{2\pi}{\log|s|}\sum_{j\ne i}(G_j-G_i)\,(\beta\cdot C_j)_p .
\]
Multiply by $F_i$ and sum over the branches. By
\eqref{eq:local-intersection} the result is
\[
 \frac{2\pi}{\log|s|}\sum_{i}\sum_{j\ne i}F_i(G_j-G_i)(C_i\cdot C_j)_p
 =-\frac{2\pi}{\log|s|}\sum_{i<j}(C_i\cdot C_j)_p(F_i-F_j)(G_i-G_j),
\]
which is \eqref{eq:flux}. Since $h_G$ is harmonic, Stokes' theorem gives
\[
 E_M(h_F,h_G)=\int_Mdh_F\wedge\dc h_G
 =\sum_\beta\int_{\gamma_\beta(r,s)}h_F\,\dc h_G .
\]
On $\gamma_\beta(r,s)$ we have $|h_F-F_{i(\beta)}|\le c(r)|F|/\log(1/|s|)$
and $|\dc h_G|\le c(r)|G|/\log(1/|s|)$ by Lemma~\ref{lem:collar}, and the
lengths of these circles are bounded. Replacing $h_F$ by $F_{i(\beta)}$
and applying \eqref{eq:flux} proves \eqref{eq:energy}.
\end{proof}

\section{The upper bound}\label{sec:upper}

\begin{proposition}\label{prop:upper}
For $1\le k\le N-1$,
$\limsup_{s\to0}\log(1/|s|)\,\lambda_k(s)\le2\pi\mu_k$.
\end{proposition}
\begin{proof}
Fix $r\in(0,\rho_0]$ and write $L=\log(1/|s|)$. For each $p\in\Sigma$ let
$\chi_p$ be a smooth function on $X$ with values in $[0,1]$, equal to one
on $B_p(r)$ and supported in $B_p(2r)$. For $F\in\RR^N$ define a function
$U_F$ on $X_s$ by
\[
 U_F=\begin{cases}
 F_i&\text{on }\Omega_i(2r,s),\\
 F_{i(\beta)}+\chi_p\,(h_F-F_{i(\beta)})&\text{on }X_s\cap W_\beta\cap
     \overline{B_p(2r)},\\
 h_F&\text{on }M_p(r,s),
 \end{cases}
\]
where $h_F$ is the function \eqref{eq:hF} at $p$. By
Lemmas~\ref{lem:thick} and~\ref{lem:collar} these prescriptions agree on
overlaps for small $s$ and define a smooth function, linear in $F$. The
energy of $U_F$ vanishes on the $\Omega_i(2r,s)$ and equals
$E_{M_p(r,s)}(h_F)$ on $M_p(r,s)$. On the collars
$X_s\cap(B_p(2r)\setminus B_p(r))$, which have bounded area,
Lemma~\ref{lem:collar} gives
$|dU_F|\le|d\chi_p||h_F-F_{i(\beta)}|+|dh_F|\le c(r)|F|/L$. Hence by
Lemma~\ref{lem:energy}
\begin{equation}\label{eq:upper-energy}
 \Bigl|E(U_F)-\frac{2\pi}{L}\,q(F)\Bigr|\le\frac{c(r)|F|^2}{L^2}.
\end{equation}
By \eqref{eq:hF-bounds}, $U_F$ is everywhere a convex combination of the
$F_i$, so $|U_F|\le|F|$. Lemma~\ref{lem:area} therefore
gives
\begin{equation}\label{eq:upper-mass}
 \Bigl|\,\|U_F\|^2-\sum_iF_i^2\area\Omega_i(2r,s)\Bigr|\le Cr^2|F|^2,
 \qquad
 \liminf_{s\to0}\|U_F\|^2\ge m(F)-Cr^2|F|^2 .
\end{equation}
Let $V_k\subset\RR^N$ be the span of eigenvectors of $q$ relative to $m$
for $\mu_0,\ldots,\mu_k$. Note that $|F|^2\le m(F)/\min_iA_i$, and that the convergence of
the areas in \eqref{eq:upper-mass} is uniform on the unit sphere of
$\RR^N$. If $r$ is small, \eqref{eq:upper-mass} shows
that $F\mapsto U_F$ is injective on $V_k$ for small $s$, so
$\{U_F:F\in V_k\}$ has dimension $k+1$. On $V_k$ we have
$q(F)\le\mu_km(F)$. The min--max principle,
\eqref{eq:upper-energy} and \eqref{eq:upper-mass} give
\[
 \limsup_{s\to0}L\,\lambda_k(s)
 \le\sup_{F\in V_k\setminus\{0\}}\frac{2\pi\,q(F)}{m(F)-Cr^2|F|^2}
 \le\frac{2\pi\mu_k}{1-Cr^2/\min_iA_i}.
\]
Let $r\to0$.
\end{proof}

\section{The lower bound}\label{sec:lower}

Throughout this section $u$ denotes an eigenfunction of $\Delta_s$ with
eigenvalue $\lambda\le1$ and $\|u\|=1$, and $r\in(0,\rho_0]$ is fixed.
Let $a_i=a_i(u)$ be the mean value of $u$ over $\Omega_i(r/2,s)$.

\begin{lemma}\label{lem:thick-estimates}
For sufficiently small $s$,
\begin{enumerate}
\item $|a_i|\le C$;
\item $\|u-a_i\|_{L^2(\Omega_i(r/2,s))}\le c(r)\lambda^{1/2}$;
\item $\sup_{\gamma_\beta(r,s)}|u-a_{i(\beta)}|\le c(r)\lambda^{1/2}$ for
every branch $\beta$ at every $p\in\Sigma$.
\end{enumerate}
\end{lemma}
\begin{proof}
By the Cauchy--Schwarz inequality,
$|a_i|\le\area(\Omega_i(r/2,s))^{-1/2}$, and (1) follows from
\eqref{eq:area-limit} once $\rho_0$ is small. For (2) and (3) transport
$u$ to the fixed surface $\Omega_i^0(r/2)$ by the diffeomorphisms of
Lemma~\ref{lem:thick}. The transported metrics converge smoothly to
$g_0$, so the Poincar\'e inequality on the connected surface
$\Omega^0_i(r/2)$ holds for them with a uniform constant; since
$E(u)=\lambda$, this is (2). The function $v=u-a_i$ satisfies
$\Delta_sv=\lambda u$, whose $L^2$ norm is at most $\lambda$. The curves
$\gamma_\beta(r,s)$ with $i(\beta)=i$ remain in a fixed compact subset of
the interior of $\Omega_i^0(r/2)$ after transport. Interior elliptic
estimates with uniformly controlled coefficients and the Sobolev embedding
$W^{2,2}\subset C^0$ in dimension two give
$\sup_{\gamma_\beta}|v|\le c(r)(\|\Delta_sv\|_{L^2}+\|v\|_{L^2})
\le c(r)\lambda^{1/2}$, the norms being taken over $\Omega_i(r/2,s)$.
\end{proof}

\begin{lemma}\label{lem:milnor-mass}
For sufficiently small $s$ and every $p\in\Sigma$,
\[
 \int_{M_p(r,s)}u^2\,dA\le Cr^2+c(r)\lambda^{1/2}.
\]
\end{lemma}
\begin{proof}
Let $M=M_p(r,s)$ and let $h$ be the harmonic function on $M$ with
boundary values $u$. Then $w=u-h$ vanishes on $\partial M$, and
$E_M(w)\le E_M(u)\le\lambda$, because $E_M(u)=E_M(h)+E_M(w)$ for $h$
harmonic and $w$ vanishing on the boundary. By the maximum
principle and Lemma~\ref{lem:thick-estimates},
$|h|\le C+c(r)\lambda^{1/2}$, so Lemma~\ref{lem:area} gives
$\int_Mh^2\,dA\le Cr^2(C+c(r)\lambda^{1/2})^2$.

For $w$ we compare with the flat metric of the chart. Let $dA_0$ be the
area form on $X_s\cap U_p$ induced by that metric. Both induced metrics
are Hermitian on the Riemann surface $X_s\cap U_p$, hence conformal; the
Dirichlet energy is the same for both, and $dA\le C\,dA_0$. Since
$d\dc|z|^2$ restricts to $4\,dA_0$ on a complex curve, Stokes' theorem
gives
\[
 4\int_Mw^2\,dA_0=-2\int_Mw\,dw\wedge\dc|z|^2
 \le4r\Bigl(\int_Mw^2\,dA_0\Bigr)^{1/2}E_M(w)^{1/2},
\]
because $|d|z|^2|\le2|z|\le2r$ on $M$. Hence
$\int_Mw^2\,dA\le Cr^2\lambda$. The lemma follows from
$u^2\le2h^2+2w^2$.
\end{proof}

Lemmas~\ref{lem:thick-estimates} and~\ref{lem:milnor-mass} extend to
linear combinations $u=\sum_{l=1}^n\zeta_lu_l$ of orthonormal eigenfunctions
with eigenvalues at most $\lambda\le1$ and $\sum_l\zeta_l^2=1$, with constants
multiplied by a factor depending only on $n$: the means $a_i(u)$ are
linear in $u$, the estimates of Lemma~\ref{lem:thick-estimates} transfer
by the triangle inequality, and $E(u)=\sum_l\zeta_l^2\lambda_l\le\lambda$, which
is all that the proof of Lemma~\ref{lem:milnor-mass} uses besides them.

\begin{lemma}[Energy lower bound]\label{lem:capacity}
For sufficiently small $s$, every $p\in\Sigma$ and $a=(a_i(u))_i$,
\[
 E_{M_p(r,s)}(u)\ge\frac{2\pi}{\log(1/|s|)}\,q_p(a)
 -\frac{c(r)}{\log(1/|s|)}\Bigl(\lambda^{1/2}+\frac1{\log(1/|s|)}\Bigr).
\]
The same inequality holds, with $c(r)$ multiplied by a constant depending
only on $n$, for $u=\sum_{l=1}^n\zeta_lu_l$, where $u_1,\ldots,u_n$ are
orthonormal eigenfunctions with eigenvalues at most $\lambda$,
$\sum_l\zeta_l^2=1$, and $a=\sum_l\zeta_la(u_l)$.
\end{lemma}
\begin{proof}
Write $M=M_p(r,s)$ and $L=\log(1/|s|)$, and let $h_a$ be the function
\eqref{eq:hF} for $F=a$. It is harmonic and smooth up to $\partial M$, so
\[
 I:=\int_Mdu\wedge\dc h_a
 =\sum_\beta a_{i(\beta)}\int_{\gamma_\beta(r,s)}\dc h_a
 +\sum_\beta\int_{\gamma_\beta(r,s)}(u-a_{i(\beta)})\,\dc h_a .
\]
The first sum equals $P:=(2\pi/L)\,q_p(a)\ge0$ by \eqref{eq:flux}. By
Lemmas~\ref{lem:collar} and~\ref{lem:thick-estimates}, the second is
bounded in absolute value by $e:=c(r)\lambda^{1/2}/L$. By
Lemma~\ref{lem:energy}, $E_M(h_a)\le P+\theta$ with
$\theta=c(r)/L^2$. If $P\le e$ there is nothing to prove. Otherwise
$I\ge P-e>0$, and the Cauchy--Schwarz inequality
$I^2\le E_M(u)E_M(h_a)$ gives
\[
 E_M(u)\ge\frac{(P-e)^2}{P+\theta}\ge P-2e-\theta .
\]
For a combination $u=\sum\zeta_lu_l$ the same argument applies, by the
paragraph preceding the lemma.
\end{proof}

\begin{proposition}\label{prop:lower}
For $1\le k\le N-1$,
\[
 \liminf_{s\to0}\log(1/|s|)\,\lambda_k(s)\ge2\pi\mu_k .
\]
Moreover
$\liminf_{s\to0}\lambda_N(s)>0$.
\end{proposition}
\begin{proof}
Let $n\le N+1$ and let $u_1,\ldots,u_n$ be orthonormal eigenfunctions of
$\Delta_s$ with eigenvalues at most $\lambda\le1$. Let $V$ be their span
and $\Theta:V\to\RR^N$ the linear map $u\mapsto(a_i(u))_i$. For
$u=\sum \zeta_lu_l$ with $\sum \zeta_l^2=1$, the triangle inequality in $L^2$ and
Lemmas~\ref{lem:thick-estimates} and~\ref{lem:milnor-mass} give
\[
 \int_{\Omega_i(r,s)}u^2\,dA\le a_i^2\area\Omega_i(r,s)
 +c(r)\lambda^{1/2},\qquad
 \int_{M_p(r,s)}u^2\,dA\le C r^2+c(r)\lambda^{1/2},
\]
with constants depending on $N$. Since $X_s$ is the union of the
$\Omega_i(r,s)$ and the $M_p(r,s)$, summing and using
\eqref{eq:area-limit} yields
\begin{equation}\label{eq:mass-lower}
 m(\Theta u)\ge\bigl(1-Cr^2-c(r)\lambda^{1/2}-o(1)\bigr)\|u\|^2
 \qquad(u\in V),
\end{equation}
where $o(1)\to0$ as $s\to0$ for fixed $r$.

\emph{The gap.} Suppose $\lambda_N(s_l)\to0$ along a sequence
$s_l\to0$. Apply the preceding with $n=N+1$, the eigenfunctions for
$\lambda_0,\ldots,\lambda_N$ and $\lambda=\lambda_N(s_l)$. Choosing $r$
with $Cr^2<1/2$, \eqref{eq:mass-lower} shows that $\Theta$ is injective for
large $l$, which is impossible since $\dim V=N+1$.

\emph{The limits.} Let $1\le k\le N-1$, $n=k+1$, and take the
eigenfunctions for $\lambda_0,\ldots,\lambda_k$, with
$\lambda=\lambda_k(s)$. By Proposition~\ref{prop:upper},
$\lambda\le C/L$ with $L=\log(1/|s|)$. For $r$ small and then $s$ small,
\eqref{eq:mass-lower} shows that $\Theta$ is injective, so $\Theta(V)$ has
dimension $k+1$ and contains a vector $\Theta u$, $\|u\|=1$, with
$q(\Theta u)\ge\mu_k\,m(\Theta u)$. For this $u$,
Lemma~\ref{lem:capacity}, summed over $p\in\Sigma$, and
\eqref{eq:mass-lower} give
\[
 \lambda_k(s)\ge E(u)\ge\sum_{p\in\Sigma}E_{M_p(r,s)}(u)
 \ge\frac{2\pi\mu_k}{L}\bigl(1-Cr^2-o(1)\bigr)
 -\frac{c(r)}{L^{3/2}} .
\]
Multiply by $L$, let $s\to0$ and then $r\to0$.
\end{proof}

Propositions~\ref{prop:upper} and~\ref{prop:lower} prove
Theorem~\ref{thm:main}.

\begin{proof}[Proof of Corollaries~\ref{cor:two} and~\ref{cor:product}]
Let $Q$ be the matrix of $q$ and $A=\operatorname{diag}(A_i)$. The
$\mu_k$ are the roots of $\det(Q-tA)$. This polynomial has leading
coefficient $(-1)^N\prod_iA_i$, no constant term, and linear coefficient
$-\sum_iA_i\det Q^{(i)}$, where $Q^{(i)}$ is obtained by deleting the
$i$th row and column. Since $Q$ is a weighted graph Laplacian, all
$\det Q^{(i)}$ are equal to the weighted number $\kappa$ of spanning trees
\cite{chaiken}. Hence
$\prod_{k\ge1}\mu_k=\kappa\sum_iA_i/\prod_iA_i$, and
Corollary~\ref{cor:product} follows from \eqref{eq:main}. For $N=2$,
$\kappa=C_1\cdot C_2$.
\end{proof}

\section{Examples and remarks}\label{sec:examples}

\begin{example}[Nodal fibres]\label{ex:nodal}
If $X_0$ has only ordinary nodes, $C_i\cdot C_j$ is the number of nodes
joining $C_i$ and $C_j$, and nodes of a single component do not enter.
The matrix of $q$ is the Laplacian of the dual graph of $X_0$ with loops
removed.
\end{example}

\begin{example}[Tangency and higher contact]\label{ex:tacnode}
Let two components be given near $p$ by $f_1=y$ and $f_2=y-x^n$, with
$\pi=f_1f_2$. Then $(C_1\cdot C_2)_p=n$: the point contributes as $n$
nodes. We sketch a direct argument. On the fibre $\{y(y-x^n)=s\}$ the
function $y$ omits zero, and each of its values other than the two square
roots of $s$ is taken at the $n$ solutions of $x^n=y-s/y$. After removing from the
local fibre and from the annulus sets on which $\log|y|/\log|s|$ has
energy $O(\log(1/|s|)^{-2})$, $y$ therefore exhibits the local fibre as an
$n$-sheeted branched cover of an annulus $\{c^{-1}|s|<|y|<c\}$. The
function $\log|y|/\log|s|$ has energy $2\pi\log(c^2/|s|)/\log(1/|s|)^2$
on the annulus and $n$ times as much on the cover.
\end{example}

\begin{example}[Ordinary multiple points]\label{ex:multiple}
Suppose $d$ distinct components pass through $p$ with pairwise distinct
tangents. Then $q_p$ is the Laplacian form of the complete graph on these
$d$ components. This follows from $(C_i\cdot C_j)_p=1$. It can be
interpreted as follows, although the theorem does not apply to the
non-reduced fibre of the blown-up family: in the blow-up of $p$ the exceptional curve
appears with multiplicity $d$ and meets each strict transform once; the
corresponding star-shaped network, with conductance $d$ on each edge and
no mass at the centre, has the complete graph with unit conductances as
its star--mesh reduction to the outer vertices, that is, the Schur
complement of its Laplacian with respect to the centre.
\end{example}

\begin{example}[The Fermat pencil]\label{ex:fermat}
Let $d\ge2$ and
\[
 X=\{(s,[x:y:z])\in\mathbb D\times\PP^2:\ x^d+y^d+sz^d=0\},
\]
with $\pi$ the first projection and $g$ the restriction of the product of
the flat metric on $\mathbb D$ and the Fubini--Study metric, so that $g_s$
is the Fubini--Study metric on the plane curve $X_s$. This is
\cite[Example~8.1]{dai-yoshikawa}. The surface $X$ is smooth, and $X_0$
consists of $d$ lines through $[0:0:1]$, with $\pi=-\prod_\zeta(x-\zeta y)$
in the coordinates $(x,y)$ at that point, the product over the $d$th
roots of $-1$. Thus $N=d$, $C_i\cdot C_j=1$ for $i\ne j$, and all $A_i$
equal the area $A$ of a line. The Laplacian of the complete graph has
the single nonzero eigenvalue $d$, of multiplicity $d-1$, so
\[
 \lambda_k(s)=\frac{2\pi d}{A}\,\frac{1+o(1)}{\log(1/|s|)}
 \qquad(1\le k\le d-1).
\]
\end{example}

\begin{remark}[Burger's theorem]\label{rem:burger}
The hyperbolic metrics on the fibres of a degenerating family of stable
curves are not induced by a smooth metric on the total space, and
Theorem~\ref{thm:main} does not apply to them. Nevertheless, after the
following conversion Burger's graph coincides with the graph of
Theorem~\ref{thm:main}. At a node with plumbing parameter $t$, that is,
with local equation $xy=t$ in the family, the short
geodesic has length $\ell=2\pi^2/\log(1/|t|)+O(\log(1/|t|)^{-2})$
\cite[Example~4.3]{wolpert}, and the hyperbolic area of a component is
$2\pi(2\gamma_v-2+n_v)$. Burger's ratio $1/(2\pi^2)$ for edge weight $\ell$ and
vertex mass $2\gamma_v-2+n_v$ \cite[Theorem~1.1]{burger} therefore corresponds
to conductance $2\pi/\log(1/|t|)$ per node and mass equal to the area, as
in Example~\ref{ex:nodal}. Burger's theorem moreover allows the lengths to
tend to zero at independent rates, which is outside the one-parameter
setting considered here.
\end{remark}

\begin{remark}[The pairing of Cao]\label{rem:cao}
Let $\mathsf M=(C_i\cdot C_j)$ be the intersection matrix, so that
$-\mathsf M$ is the matrix of $q$. For a K\"ahler metric on $X$ and reduced $X_0$, Cao
considers the pairing
$\langle\alpha,\alpha'\rangle(s)=\int_{X_s}\varphi_s\alpha'$ of smooth
two-forms with vanishing fibre integrals, where $\varphi_s$ solves
$dd^c\varphi_s=-\alpha|_{X_s}$ with his normalization $d^c=\frac{i}{4\pi}(\bar\partial-\partial)$, and proves
that $\langle\alpha,\alpha'\rangle-c_{\alpha,\alpha'}\log|s|^2$ extends
continuously to $\mathbb D$, with $c_{\alpha,\alpha'}=v_\alpha^{T}\mathsf M^+v_{\alpha'}$,
$v_\alpha=(\int_{C_i}\alpha)_i$, and $\mathsf M^+$ the Moore--Penrose inverse
\cite[Theorem~0.2]{cao-pairing}. Since
$\langle\alpha,\alpha\rangle$ equals
$4\pi\sum_{k\ge1}\lambda_k^{-1}(\int_{X_s}\Phi_k\alpha)^2$ for an orthonormal
basis of eigenfunctions $\Phi_k$, the logarithmic term is carried by the
$N-1$ small eigenvalues. Cao also shows that the corresponding
eigenfunctions are, up to errors of order $\log(1/|s|)^{-1/2}$ in $L^2$,
combinations of functions that are constant on the parts of $X_s$ near the
components \cite[Theorem~2.10]{cao-pairing}. We sketch how the limits \eqref{eq:main} follow. The period vectors
$v_\alpha$ exhaust the vectors with zero sum, and from these statements
and the two-sided bounds of \cite{dai-yoshikawa} one can deduce, by
comparing two quadratic forms on that space along subsequences, that
$\lambda_k(s)\log(1/|s|)$ converges to $2\pi\mu_k$. Thus for K\"ahler
metrics the limits \eqref{eq:main} also follow from \cite{cao-pairing}
together with \cite{dai-yoshikawa}, although they are not stated there;
the two results describe the same leading behaviour of the Green operator
on the small eigenspaces from complementary sides.
\end{remark}

\begin{remark}[A heuristic of Ji and Wentworth]\label{rem:jw}
For a plumbing family with a separating node and the Bergman metric,
which is not induced from the total space, \cite[Remark~5.10]{ji-wentworth}
suggests that $\lambda_1$ should behave like the first Dirichlet
eigenvalue of one of the two components with a disc of radius
$\varepsilon\sim|t|^{1/2}$ removed, namely
$2\pi/(A_i\log(1/\varepsilon))$ to leading order, and concludes that
$\lambda_1$ is expected to be bounded above by a multiple of
$1/\log(1/|t|)$. Taken literally the heuristic gives
$4\pi/(A_i\log(1/|t|))$. For induced metrics, Corollary~\ref{cor:two}
gives $2\pi(1/A_1+1/A_2)/\log(1/|t|)$ at a single node, which coincides
with that value exactly when $A_1=A_2$.
\end{remark}

\begin{remark}[Hypotheses]\label{rem:hypotheses}
The Hermitian condition on $g$ is used only in the balls $B_p(4\rho_0)$,
where it makes the Dirichlet energy on the fibres conformally invariant
and the functions $h_F$ harmonic. If $X_0$ is not reduced the functions
\eqref{eq:hF} are still defined, with $f_i$ replaced by $f_i^{m_i}$ for the
multiplicities $m_i$, but the fibre near a multiple component is a
covering of it that need not be connected, and according to \cite[Theorem~0.4]{cao-kahler} the number of small
eigenvalues is the number of components of the central fibre after a base
change and normalization, minus one; we do not treat
this case. In higher dimension, for
two-forms on Calabi--Yau threefolds in conifold transitions with
Ricci-flat metrics, it is shown in the preprint \cite{wuebben-conifold}
that the small eigenvalues have the same rate and the same
constant $2\pi$ per node, with the intersection form on divisor classes in
the role of $m$ and the periods over the exceptional curves in the role of
the differences $F_i-F_j$.
\end{remark}

\begin{thebibliography}{BHPV04}

\bibitem[BHPV04]{bhpv}
W.~P. Barth, K. Hulek, C.~A.~M. Peters and A. Van de Ven,
\emph{Compact complex surfaces}, 2nd ed.,
Ergebnisse der Mathematik und ihrer Grenzgebiete (3), vol.~4,
Springer, Berlin, 2004.

\bibitem[Bur90]{burger}
M. Burger,
\emph{Small eigenvalues of Riemann surfaces and graphs},
Math. Z. \textbf{205} (1990), 395--420.
\href{https://doi.org/10.1007/BF02571252}{doi:10.1007/BF02571252}.

\bibitem[Cao26a]{cao-pairing}
J. Cao,
\emph{The Archimedean height pairing for differential forms on degeneration of
Riemann surfaces},
preprint, 2026.
\href{https://arxiv.org/abs/2603.04802v1}{arXiv:2603.04802v1}.

\bibitem[Cao26b]{cao-kahler}
J. Cao,
\emph{Asymptotics of small eigenvalues on degenerations of K\"ahler manifolds},
preprint, 2026.
\href{https://arxiv.org/abs/2605.08023v1}{arXiv:2605.08023v1}.

\bibitem[Cha82]{chaiken}
S. Chaiken,
\emph{A combinatorial proof of the all minors matrix tree theorem},
SIAM J. Algebraic Discrete Methods \textbf{3} (1982), no.~3, 319--329.
\href{https://doi.org/10.1137/0603033}{doi:10.1137/0603033}.

\bibitem[DY26]{dai-yoshikawa}
X. Dai and K.-I. Yoshikawa,
\emph{Degeneration of Riemann surfaces and small eigenvalues of the Laplacian},
to appear in Courant Journal of Pure and Applied Mathematics, 2026 preprint.
\href{https://arxiv.org/abs/2509.06151v3}{arXiv:2509.06151v3}.

\bibitem[JW92]{ji-wentworth}
L. Ji and R. Wentworth,
\emph{Spectral convergence on degenerating surfaces},
Duke Math. J. \textbf{66} (1992), no.~3, 469--501.
\href{https://doi.org/10.1215/S0012-7094-92-06615-4}{doi:10.1215/S0012-7094-92-06615-4}.
Correction: Duke Math. J. \textbf{90} (1997), no.~1, 205--207.

\bibitem[Mil68]{milnor}
J. Milnor,
\emph{Singular points of complex hypersurfaces},
Annals of Mathematics Studies, no.~61,
Princeton University Press, Princeton, NJ, 1968.

\bibitem[SWY80]{swy}
R. Schoen, S. Wolpert and S.-T. Yau,
\emph{Geometric bounds on the low eigenvalues of a compact surface},
in \emph{Geometry of the Laplace operator},
Proc. Sympos. Pure Math., vol.~36,
American Mathematical Society, Providence, RI, 1980, 279--285.
\href{https://doi.org/10.1090/pspum/036/573440}{doi:10.1090/pspum/036/573440}.

\bibitem[Wol90]{wolpert}
S.~A. Wolpert,
\emph{The hyperbolic metric and the geometry of the universal curve},
J. Differential Geom. \textbf{31} (1990), no.~2, 417--472.
\href{https://doi.org/10.4310/jdg/1214444322}{doi:10.4310/jdg/1214444322}.

\bibitem[W26]{wuebben-conifold}
B.~J.~Wuebben,
\emph{Integral topology and small eigenvalues in Calabi--Yau conifold
transitions},
preprint, November~9, 2026.
\href{https://doi.org/10.5281/zenodo.24096887}{doi:10.5281/zenodo.24096887}.

\end{thebibliography}
\end{document}
